\chapter{Introduction}
\label{ch:intro}

% da scrivere per ultima; la descrizione di Polymarket non va qui ma nel background
\todo{to be written last, once the other chapters are stable.  Only the outline for now.}

\section{Context}
\label{sec:intro:object}

Prediction markets as aggregators of dispersed information
\cite{wolfers2004prediction,arrow2008promise}; a decentralised one, Polymarket, exposes
the whole order book and every trade with the wallet behind it, at a resolution of about
two seconds, on tens of thousands of markets at once.  \todo{one page on why this is a
different object for a data scientist; the mechanics stay in \cref{ch:pm}}

\section{The problem}
\label{sec:intro:difficulty}

\todo{the task (direction of the mid price over one minute to twelve hours) and why it is
hard: a price consolidates information; it is bounded and discrete; it is mostly still;
much of what moves it is outside the data}

\section{Research questions}
\label{sec:intro:questions}

\begin{enumerate}[nosep,label=\textbf{Q\arabic*},ref=Q\arabic*]
  \item \label{q:process} What regularities does the price process of a prediction
  market have, and how is its trading organised?
  \item \label{q:information} Which publicly observable information predicts the
  direction of the mid price?
  \item \label{q:models} Which families of models can read that information, and what
  is a sequential state or a neighbourhood worth over a stateless learner on the same
  inputs?
  \item \label{q:horizon} How do the answers change with the horizon and with the kind
  of market?
\end{enumerate}

\section{Contributions}
\label{sec:intro:contributions}

\todo{list, once the results chapter is in place}

\section{Structure of the thesis}
\label{sec:intro:structure}

\todo{one paragraph per chapter}
